\documentclass[11pt]{article}

\usepackage[margin=1in]{geometry}
\usepackage{amsmath,amssymb,amsthm,mathtools}
\usepackage{graphicx}
\usepackage{microtype}
\usepackage{booktabs,tabularx}
\usepackage{enumitem}
\usepackage[dvipsnames]{xcolor}
\usepackage[colorlinks=true,linkcolor=MidnightBlue,citecolor=MidnightBlue,urlcolor=MidnightBlue]{hyperref}

\newtheorem{theorem}{Theorem}
\newtheorem{lemma}[theorem]{Lemma}
\newtheorem{corollary}[theorem]{Corollary}
\newtheorem{remark}[theorem]{Remark}

\newcommand{\X}{\mathcal X}
\newcommand{\Cb}{C_b}
\newcommand{\ch}{\operatorname{ch}}
\newcommand{\ran}{\operatorname{ran}}
\newcommand{\spec}{\sigma}

\setlist{itemsep=3pt,topsep=5pt}

\title{All Six Open Questions on the Differentiation Operator\\
\large A norm counterexample and universal power, range, and compactness obstructions}
\author{Automated mathematical discovery run \texttt{fa\_banach\_001}\\
Agent \texttt{agent\_lane\_14} \quad Model GPT5.6}
\date{22 July 2026}

\begin{document}
\maketitle

\begin{center}
\fcolorbox{MidnightBlue}{MidnightBlue!5}{%
\parbox{0.91\textwidth}{\textbf{Status: full solution of all six questions,
likely valid, pending human review.} The packet answers every question in
Section 6 of Allen--Jackson \cite{AllenJackson}, including exact spectra on
the weighted Banach and discrete Hardy spaces.}}
\end{center}

\section{Result at a glance}

Let $T$ be the locally finite rooted tree of the source paper, with every
vertex having a child, let $b(o)=o$ and $b(v)=v^-$ for $v\ne o$, and put
\[
  \Cb f=f\circ b,\qquad Df=f-f\circ b.
\]
On a discrete function space, point evaluations are continuous and the
functions separate every pair of vertices.

\begin{center}
\begin{tabularx}{0.96\textwidth}{@{}c>{\bfseries}lX@{}}
\toprule
Question & Answer & Mechanism \\
\midrule
1 & Yes & A weighted Hilbert space has $\|\Cb\|=\sqrt2$ but $\|D\|=2$. \\
2 & No & Pointwise bounded $D$-orbits contain only constant functions; hence $D$ is never power-bounded. \\
3 & No & Every function in $\ran(\Cb)$ agrees at the root and each root child. \\
4 & No & A nonzero spectral point of $D$ is not an eigenvalue, ruling out compactness. \\
5 & Exact formulas & A growth-radius disk plus the root point for weighted spaces; the unit disk for Hardy spaces. \\
6 & No & The compactness obstruction is universal, hence applies to the intended path-tree spaces. \\
\bottomrule
\end{tabularx}
\end{center}

Thus the norm conjecture in the source is false, while its isometry conjecture
is true in the stronger power-bounded form.

\section{Source questions}

The questions occur in Section 6, page 17 of arXiv:2207.10211
\cite{AllenJackson}. The following is a full-width crop of the source page.

\begin{figure}[ht]
  \centering
  \includegraphics[width=\textwidth]{figures/open_problem_crop.png}
  \caption{The six open questions in the source paper. This packet answers
  all six.}
\end{figure}

The relevant questions ask: whether $\|D\|$ can differ from
$1+\|\Cb\|$; whether $D$ can act isometrically; whether $\Cb$ can be
surjective; whether $D$ can be compact; and whether $D$ is compact on the
Lipschitz or weighted Banach space of a path tree.

\section{Definitions and conventions}

A \emph{discrete function space} $\X$ on $T$ is a complex Banach space of
complex-valued functions on $T$ such that:
\begin{enumerate}
  \item every evaluation $\delta_v(f)=f(v)$ is continuous;
  \item if $f(v)=f(w)$ for every $f\in\X$, then $v=w$.
\end{enumerate}
The second condition is the point-separation axiom. The source assumes that
every vertex has a child. This child condition, rather than homogeneity or a
specific branching number, is exactly what the proofs below use.

Whenever $D$ maps $\X$ into itself, the source's closed graph argument makes
$D$ bounded, and $\Cb=I-D$ is bounded as well \cite[Theorems 2.6--2.7]{AllenJackson}.

\section{Proof intuition}

Four elementary geometric features drive all six answers.

First, distribute one unit of Hilbert mass from every parent among its
children. After conjugating the resulting weighted $\ell_2$ space to ordinary
$\ell_2(T)$, the operator $\Cb$ becomes $P_o+U$, where $P_o$ is the root
projection and $U$ shifts mass one level outward isometrically. The extra root
copy raises $\|\Cb\|$ to $\sqrt2$, while alternating long $U$-orbits show that
$\|I-P_o-U\|=2$. This breaks the conjectured norm identity.

Second, repeated differences reveal any nonzero edge increment one generation
later as a term growing linearly with the iteration number. Thus a function
with a pointwise bounded $D$-orbit must be constant. Norm power-boundedness
would make every orbit pointwise bounded, which is incompatible with point
separation.

Third, the root and each root child lie in the same fibre of $b$. Hence every
function in the range of $\Cb$ identifies those two vertices, so $\Cb$ cannot
be onto a point-separating space. This range obstruction puts $1$ in the
spectrum of $D$ but not in its point spectrum, which compact operator theory
forbids.

Fourth, both named spaces split as $\ell_\infty$-sums of their finite
dimensional level spaces. On the nonroot tail, $\Cb$ is an operator-weighted
unilateral shift. A one-level resolvent test shows that the spectrum of this
shift is exactly the disk determined by its uniform asymptotic block-growth
rate. For $L_\mu^\infty$ this rate is a descendant/ancestor weight ratio; for
$\mathbb T_p$ every level map is isometric, so the rate is one.

\section{A Hilbert-space counterexample to the norm identity}

\begin{theorem}[Answer to Question 1]\label{thm:norm}
On every rooted locally finite tree in which each vertex has a child, there is
a discrete Hilbert function space $H_w$ such that
\[
  \|\Cb\|=\sqrt2,\qquad \|D\|=2.
\]
Consequently $\|D\|\ne 1+\|\Cb\|$.
\end{theorem}

\begin{proof}
For every nonroot vertex $v$, choose $p_v>0$ so that
\[
  \sum_{v\in\ch(u)}p_v=1
  \qquad (u\in T).
\]
This is possible because every child set is finite and nonempty. Define
$w(o)=1$ and recursively
\[
  w(v)=p_vw(v^-),\qquad v\ne o.
\]
All weights are positive. Let
\[
  H_w=\left\{f:T\longrightarrow\mathbb C:
  \|f\|_w^2:=\sum_{v\in T}w(v)|f(v)|^2<\infty\right\}.
\]
This is a Hilbert space. Evaluation at $v$ has norm $w(v)^{-1/2}$, and the
finitely supported point masses belong to $H_w$ and separate the vertices.
Thus $H_w$ is a discrete function space.

The map $J:H_w\to\ell_2(T)$ given by
\[
  (Jf)(v)=w(v)^{1/2}f(v)
\]
is unitary. Define $U$ on $\ell_2(T)$ by
\[
  (Uy)(o)=0,\qquad (Uy)(v)=p_v^{1/2}y(v^-)\quad(v\ne o).
\]
Then
\[
\begin{aligned}
  \|Uy\|_2^2
  &=\sum_{u\in T}\sum_{v\in\ch(u)}p_v|y(u)|^2\\
  &=\sum_{u\in T}|y(u)|^2,
\end{aligned}
\]
so $U$ is an isometry. Let $P_o$ be the orthogonal projection onto the root
coordinate. Directly from the definitions,
\[
  J\Cb J^{-1}=P_o+U,
  \qquad JDJ^{-1}=I-P_o-U=(I-P_o)(I-U).
\]
The vectors $P_oy$ and $Uy$ have disjoint supports. Therefore
\[
  \|(P_o+U)y\|_2^2=\|P_oy\|_2^2+\|y\|_2^2
  \le 2\|y\|_2^2.
\]
Equality holds for the root unit vector $e_o$, so $\|\Cb\|=\sqrt2$.

The displayed factorization gives $\|D\|\le2$. To obtain the reverse
inequality, observe that $U^ke_o$ is a unit vector supported on level $k$;
hence these vectors are orthonormal. For $N\ge1$, put
\[
  y_N=\frac1{\sqrt N}\sum_{k=0}^{N-1}(-1)^kU^ke_o.
\]
Then $\|y_N\|_2=1$. After applying $(I-P_o)(I-U)$, the coefficients on
levels $1,\ldots,N-1$ have modulus $2/\sqrt N$, and the final coefficient on
level $N$ has modulus $1/\sqrt N$. Hence
\[
  \|Dy_N\|_2^2=\frac{4(N-1)+1}{N}=4-\frac3N.
\]
Letting $N\to\infty$ yields $\|D\|\ge2$. Thus $\|D\|=2$, whereas
$1+\|\Cb\|=1+\sqrt2$.
\end{proof}

\section{Iterated differences cannot be pointwise bounded}

\begin{lemma}[Pointwise growth lemma]\label{lem:growth}
Let $f:T\to\mathbb C$ be any function. If, for every $v\in T$, the scalar
sequence $\{(D^nf)(v):n\ge0\}$ is bounded, then $f$ is constant.
\end{lemma}

\begin{proof}
We prove by induction on $m\ge1$ that $f(v)=f(v^-)$ for every vertex $v$ of
depth $m$. Fix such a $v$ and choose a child $z$ of $v$. Along the path from
the root through $v$ to $z$, write
\[
  a_j=f(v_j)-f(v_{j-1}),
  \qquad o=v_0,v_1,\ldots,v_m=v,v_{m+1}=z.
\]
For $m=1$, $D^nf(o)=0$ for every $n\ge1$. For the induction step, the earlier
edge increments vanish by hypothesis, which again gives
$D^nf(v_{m-1})=0$ for all $n\ge1$. In either case the recurrence
\[
  (D^{n+1}f)(x)=(D^nf)(x)-(D^nf)(x^-)
\]
shows that
\[
  (D^nf)(v)=a_m\qquad(n\ge1).
\]
Applying the same recurrence at $z$ gives
\[
  (D^nf)(z)=a_{m+1}-(n-1)a_m\qquad(n\ge1).
\]
The left side is bounded in $n$, so $a_m=0$. Since $v$ was arbitrary, the
induction closes at every depth. All edge increments vanish, and connectedness
makes $f$ constant.
\end{proof}

\begin{theorem}[Answer to Question 2]\label{thm:power}
The differentiation operator is not power-bounded on any discrete function
space. In particular, it is never an isometry.
\end{theorem}

\begin{proof}
Suppose $D$ were power-bounded on $\X$, say
$\sup_n\|D^n\|\le M$. For $f\in\X$ and $v\in T$, continuity of evaluation
gives
\[
  |(D^nf)(v)|\le \|\delta_v\|M\|f\|.
\]
Lemma~\ref{lem:growth} makes every $f\in\X$ constant. Such a space cannot
separate the root from a child, contradicting the definition of a discrete
function space. An isometry has $\|D^n\|=1$ for every $n$, so it is excluded.
\end{proof}

\section{The backward composition operator is never onto}

\begin{theorem}[Answer to Question 3]\label{thm:onto}
The operator $\Cb$ is not surjective on any discrete function space.
\end{theorem}

\begin{proof}
Let $v$ be any child of the root. Every $g=\Cb f$ satisfies
\[
  g(o)=f(b(o))=f(o)=f(b(v))=g(v).
\]
If $\Cb$ were onto $\X$, every member of $\X$ would take equal values at the
distinct vertices $o$ and $v$. This contradicts point separation.
\end{proof}

This sharpens the source paper's observation that surjectivity would force
range functions to be constant on sibling sets: the fibre
$b^{-1}(o)$ already contains both the root and every root child.

\section{Differentiation is never compact}

We first close a finite-dimensional loophole needed for the compactness
argument.

\begin{lemma}\label{lem:infinite}
There is no finite-dimensional discrete function space $\X$ for which
$D\X\subseteq\X$.
\end{lemma}

\begin{proof}
The source proves that the only possible eigenvalue of $D$ is $0$, and that
the corresponding eigenvectors are constant functions
\cite[Theorem 2.14]{AllenJackson}. For completeness, a nonzero eigenvalue is
also ruled out directly by evaluating $Df=\lambda f$ first at the root and
then inductively along the tree.

If $\X$ were finite-dimensional and nonzero, complex finite-dimensional
spectral theory would therefore make $D$ nilpotent. Let $m$ be its
nilpotency index. If $m\ge2$, choose $f$ with
$g=D^{m-1}f\ne0$. Then $Dg=0$, so $g$ is constant. But $g$ lies in the range
of $D$, and every function in that range vanishes at the root. Thus $g=0$, a
contradiction. Hence $m=1$ and $D=0$. Then every member of $\X$ is constant,
again contradicting point separation.
\end{proof}

\begin{theorem}[Answers to Questions 4 and 6]\label{thm:compact}
The differentiation operator is not compact on any discrete function space.
Consequently it is not compact on either the Lipschitz space or any weighted
Banach space of the intended path tree for which it acts boundedly.
\end{theorem}

\begin{proof}
Suppose $D$ were compact on $\X$. Then $D$ is bounded, so $\Cb=I-D$ is
bounded. Lemma~\ref{lem:infinite} makes $\X$ infinite-dimensional.

By Theorem~\ref{thm:onto}, $\Cb$ is not invertible. Therefore
\[
  0\in\spec(\Cb),\qquad
  1\in\spec(D),
\]
the latter because $D-I=-\Cb$. On the other hand, $\Cb$ is injective: for
each $u\in T$, choose a child $v$ with $b(v)=u$; if $\Cb f=0$, then
$f(u)=f(b(v))=0$. Hence $1$ is not an eigenvalue of $D$.

Every nonzero spectral value of a compact operator on an infinite-dimensional
complex Banach space is an eigenvalue \cite[Chapter VII, Section 4]{Conway}.
The spectral point $1$ contradicts this theorem. Thus $D$ cannot be compact.

For a named path-tree space, an unbounded $D$ is automatically not compact;
when $D$ is bounded, the universal conclusion just proved applies.
\end{proof}

\section{Exact spectra on the weighted and Hardy spaces}

We begin with a self-contained shift lemma. It isolates the only spectral
calculation needed for both families.

\begin{lemma}[Operator-weighted unilateral shift]\label{lem:shift-spectrum}
Let $(E_n)_{n\ge1}$ be nonzero complex Banach spaces and let
$B_n:E_{n-1}\to E_n$ $(n\ge2)$ be uniformly bounded. On
\[
  Z=\left(\bigoplus_{n\ge1}E_n\right)_{\ell_\infty}
\]
define
\[
  (Sx)_1=0,\qquad (Sx)_n=B_nx_{n-1}\quad(n\ge2).
\]
For $k\ge1$, put
\[
  R_k=\sup_{n\ge1}\|B_{n+k}B_{n+k-1}\cdots B_{n+1}\|,
  \qquad
  \rho=\lim_{k\to\infty}R_k^{1/k}.
\]
Then the limit exists and
\[
  \spec(S)=\{\lambda\in\mathbb C:|\lambda|\le\rho\}.
\]
\end{lemma}

\begin{proof}
The sequence $(R_k)$ is submultiplicative, and direct calculation gives
$\|S^k\|=R_k$. Thus the limit defining $\rho$ exists and equals the spectral
radius of $S$. In particular, the spectrum is contained in the stated disk.

The operator $S$ is not onto because the first coordinate of every range
vector vanishes, so $0\in\spec(S)$. Now fix $0<|\lambda|<\rho$ and suppose,
for a contradiction, that $\lambda I-S$ has a bounded inverse of norm $C$.
Fix $n\ge1$ and a unit vector $y_n\in E_n$, and let $y\in Z$ have this as its
only nonzero coordinate. For $x=(\lambda I-S)^{-1}y$, the triangular
recurrence first gives $x_j=0$ for $j<n$, and then
\[
  x_{n+k}=\lambda^{-k-1}
  B_{n+k}B_{n+k-1}\cdots B_{n+1}y_n
  \qquad(k\ge0),
\]
where the empty product for $k=0$ is the identity. Since $\|x\|\le C$,
\[
  \|B_{n+k}\cdots B_{n+1}y_n\|
  \le C|\lambda|^{k+1}.
\]
Taking the supremum over unit $y_n$ and over $n$ yields
$R_k\le C|\lambda|^{k+1}$, hence $\rho\le|\lambda|$, a contradiction.
Therefore the open disk of radius $\rho$ lies in the spectrum. Closedness of
the spectrum supplies its boundary. (When $\rho=0$, the spectral-radius
inclusion and $0\in\spec(S)$ give the same conclusion.)
\end{proof}

\subsection{Weighted Banach spaces}

Let $\mu:T\to(0,\infty)$ be a weight and suppose
\[
  M_\mu:=\sup_{v\in T^*}\frac{\mu(v)}{\mu(v^-)}<\infty,
\]
which is exactly the boundedness condition for $D$ (equivalently $\Cb$) on
$L_\mu^\infty(T)$ in the source. For $k\ge1$, define
\[
  R_k(\mu)=\sup\left\{
  \frac{\mu(v)}{\mu(u)}:
  u\in T^*,\ v\text{ is a descendant of }u,\ |v|-|u|=k
  \right\},
\]
and
\[
  \rho_\mu=\lim_{k\to\infty}R_k(\mu)^{1/k}.
\]
The limit exists because $R_{k+\ell}(\mu)\le R_k(\mu)R_\ell(\mu)$, and
$0\le\rho_\mu\le M_\mu$.

\begin{theorem}[Weighted-space part of Question 5]\label{thm:weighted-spectrum}
If $D$ is bounded on $L_\mu^\infty(T)$, then
\[
\boxed{\quad
  \spec(\Cb)=\{1\}\cup\{\lambda:|\lambda|\le\rho_\mu\},
  \qquad
  \spec(D)=\{0\}\cup\{\lambda:|\lambda-1|\le\rho_\mu\}.
\quad}
\]
\end{theorem}

\begin{proof}
Conjugate $L_\mu^\infty(T)$ isometrically to $\ell_\infty(T)$ by
$(Jf)(v)=\mu(v)f(v)$, and group the coordinates by levels. Put
$E_n=\ell_\infty(\{v:|v|=n\})$. The map from level $n-1$ to level $n$ is
\[
  (B_nx)(v)=\frac{\mu(v)}{\mu(v^-)}x(v^-).
\]
On the nonroot tail $Z=(\bigoplus_{n\ge1}E_n)_{\ell_\infty}$, the restriction
of $J\Cb J^{-1}$ has the form in Lemma~\ref{lem:shift-spectrum}. Moreover,
\[
 \|B_{n+k}\cdots B_{n+1}\|
 =\max_{\substack{|u|=n,\ v\succeq u\\|v|-|u|=k}}
 \frac{\mu(v)}{\mu(u)},
\]
so its spectral radius is exactly $\rho_\mu$.

Splitting off the root coordinate gives the lower triangular block form
\[
  J\Cb J^{-1}=
  \begin{pmatrix}
    1&0\\ B_1&S
  \end{pmatrix}
  \quad\text{on }\mathbb C\oplus Z.
\]
For $\lambda\ne1$, elementary block elimination shows that
$J\Cb J^{-1}-\lambda I$ is invertible if and only if $S-\lambda I$ is
invertible: the scalar block $1-\lambda$ can be eliminated, leaving the
diagonal blocks $1-\lambda$ and $S-\lambda I$. For $\lambda=1$, the first
coordinate of every range vector of $J\Cb J^{-1}-I$ is zero, so the operator
is not onto. Consequently $\spec(\Cb)=\{1\}\cup\spec(S)$, and
Lemma~\ref{lem:shift-spectrum} gives the first formula. Finally, $D=I-\Cb$
and elementary spectral mapping give $\spec(D)=1-\spec(\Cb)$, which is the
second formula.
\end{proof}

The singleton in each formula matters when $\rho_\mu<1$: it records the root
block and need not be absorbed by the disk.

\subsection{Discrete Hardy spaces}

Let $T$ be $(q+1)$-homogeneous and $1\le p<\infty$. Write
\[
  c_0=1,\qquad c_n=(q+1)q^{n-1}\quad(n\ge1),
\]
and let $E_n$ be the normalized finite-dimensional level space with norm
\[
  \|x\|_{E_n}=\left(\frac1{c_n}\sum_{|v|=n}|x(v)|^p\right)^{1/p}.
\]
Then $\mathbb T_p(T)$ is isometrically
$(\bigoplus_{n\ge0}E_n)_{\ell_\infty}$.

\begin{theorem}[Hardy-space part of Question 5]\label{thm:hardy-spectrum}
For every $1\le p<\infty$,
\[
\boxed{\quad
  \spec(\Cb:\mathbb T_p\to\mathbb T_p)=\{\lambda:|\lambda|\le1\},
  \qquad
  \spec(D:\mathbb T_p\to\mathbb T_p)=\{\lambda:|\lambda-1|\le1\}.
\quad}
\]
\end{theorem}

\begin{proof}
The level map $B_n:E_{n-1}\to E_n$ replicates each parent value across its
children. For $n\ge2$, every level-$(n-1)$ vertex has $q$ children and
$c_n=qc_{n-1}$, so
\[
  \|B_nx\|_{E_n}^p
  =\frac q{c_n}\sum_{|u|=n-1}|x(u)|^p
  =\|x\|_{E_{n-1}}^p.
\]
Thus every tail map and every consecutive tail product is an isometry. In
Lemma~\ref{lem:shift-spectrum}, $R_k=1$ and $\rho=1$. Splitting off the root
and applying the same scalar-block elimination used in
Theorem~\ref{thm:weighted-spectrum} adds only the spectral point $1$, which is
already contained in the closed unit disk. Hence
$\spec(\Cb)=\{\lambda:|\lambda|\le1\}$. Translating the spectrum by
$D=I-\Cb$ yields the disk centered at $1$ of radius $1$.
\end{proof}

Theorems~\ref{thm:weighted-spectrum} and \ref{thm:hardy-spectrum} completely
answer Question 5. Together with Theorems~\ref{thm:norm}, \ref{thm:power},
\ref{thm:onto}, and \ref{thm:compact}, every question in the source list is
settled.

\section{Verification notes}

The proof was subjected to the following checks.
\begin{enumerate}
  \item The weighted space contains every point mass and all weights are
  positive, so the functional-space axiom is genuinely satisfied.
  \item The exact $U$-isometry calculation includes the root column; the root
  copy is precisely why $\|\Cb\|=\sqrt2$ rather than $1$.
  \item The lower bound for $\|D\|$ uses orthogonal level supports, not an
  unsupported approximate-eigenvector assertion.
  \item In Lemma~\ref{lem:growth}, the child one generation below the tested
  edge is essential and is supplied by the source's leafless/tree-child
  convention.
  \item The compactness proof separately excludes finite-dimensional spaces
  before invoking the compact spectral theorem.
  \item In Lemma~\ref{lem:shift-spectrum}, a one-level input forces all
  earlier solution coordinates to vanish, so the resolvent product estimate
  has no hidden homogeneous term.
  \item The weighted radius starts at nonroot vertices because the root is
  represented separately by the scalar block $1$.
  \item The Hardy replication maps are exactly isometric after normalized
  level counting, including all consecutive products used by the shift lemma.
\end{enumerate}

A reusable finite binary-tree check is included as
\texttt{code/verify\_weighted\_hilbert.py}. At depth 10 it returns
\[
 \|\Cb\|=1.414213562373,
 \quad \|D\|_{\mathrm{trunc}}=1.979642883762,
 \quad \|Dy_{10}\|^2=3.7=4-3/10.
\]
This finite computation is a consistency check only; the proofs above are
exact and infinite-dimensional.

\section{Scope, limitations, and novelty check}

The packet fully answers all six questions under the definitions and tree
conventions of the source. For the weighted-space spectrum, the formula is
stated precisely in the bounded case, since an unbounded operator has no
Banach-operator spectrum. The packet does not claim that a particular later
paper has missed these observations; it records a candidate new solution
pending expert priority review.

The bounded novelty search on 22 July 2026 covered:
\begin{itemize}
  \item the run's registry, solution, attempt, and proof-gap indexes;
  \item exact searches for arXiv:2207.10211, its title, and the source's open
  question wording;
  \item close searches combining discrete differentiation on trees with
  power-boundedness, isometry, surjectivity, and compactness;
  \item exact and close searches for the weighted Banach and Hardy spectra,
  including operator-weighted unilateral shift and $\ell_\infty$ shift
  formulations;
  \item the later citing paper arXiv:2308.11016 on forward and backward shifts
  on Hardy spaces of a tree.
\end{itemize}
No later source explicitly stating any of these six resolutions was found.
Classical operator-weighted shift literature is relevant background, but the
shift lemma used here is proved directly and the exact application to Question
5 was not located. This supports moderate novelty confidence, not an
exhaustive claim of priority.

\section*{Human review recommendation}

\textbf{Verify and circulate after a short priority search.} The decisive
checks are the exact norm calculation in Theorem~\ref{thm:norm}, the
pointwise-growth induction in Lemma~\ref{lem:growth}, and the finite-dimensional
step in Lemma~\ref{lem:infinite}. For the new spectral completion, check the
resolvent estimate in Lemma~\ref{lem:shift-spectrum} and the indexing of
$\rho_\mu$ in Theorem~\ref{thm:weighted-spectrum}.

\begin{thebibliography}{9}

\bibitem{AllenJackson}
R. F. Allen and C. M. Jackson,
\emph{The differentiation operator on discrete function spaces of a tree},
Involve \textbf{15} (2022), no. 1, 163--184;
arXiv:\href{https://arxiv.org/abs/2207.10211}{2207.10211}.

\bibitem{Conway}
J. B. Conway,
\emph{A Course in Functional Analysis}, second edition,
Graduate Texts in Mathematics 96, Springer-Verlag, New York, 1990.

\end{thebibliography}

\end{document}
